\section*{الفصل \ref{ch:b3:supramolecular} --- الكيمياء فوق الجزيئية}

\begin{solution}{exo:b3:supramolecular:1}
أيون--ثنائي قطب؛ تراص $\pi$؛ كاتيون--$\pi$؛ \omterm{def:b3:supramolecular:interactions}{رابطة هالوجينية} (من I إلى N)؛ ثلاث روابط هيدروجينية.
\end{solution}

\begin{solution}{exo:b3:supramolecular:2}
$\Delta G^\circ = -RT\ln K = -8.314 \times 298 \times \ln(\num{1.0e4}) = \qty{-22.8}{kJ/mol}$.
\end{solution}

\begin{solution}{exo:b3:supramolecular:3}
$K_{\ce{K+}}/K_{\ce{Na+}} = 10^{1.7} = 50$؛ والفرق $RT\ln 50 = \qty{9.7}{kJ/mol}$ لصالح البوتاسيوم.
\end{solution}

\begin{solution}{exo:b3:supramolecular:4}
$n/(n + m) = 1/3$: \omterm{def:b3:supramolecular:supramolecular}{مضيف} واحد لضيفين، $\mathrm{HG_2}$.
\end{solution}

\begin{solution}{exo:b3:supramolecular:5}
$H_0 + G_0 + 1/K = \qty{4.667e-3}{mol/L}$؛ $[\mathrm{HG}] = (4.667 - \sqrt{4.667^2 - 4 \times 2.0 \times 2.0})/2 \times 10^{-3} = \qty{1.13e-3}{mol/L}$: والنسبة
المرتبطة 0.566، $\delta = 3.560 + 0.160 \times 0.566 = \qty{3.651}{ppm}$.
\end{solution}

\begin{solution}{exo:b3:supramolecular:6}
$\log K = 18.3 - 8.6 = 9.7$، $K = \num{5e9}$؛ $\Delta G^\circ = -RT\ln K = \qty{-55}{kJ/mol}$. روابط Ni--N الست من
النوع نفسه في الطرفين؛ ويرتفع عدد الجزيئات الحرة من أربعة إلى
سبعة: فكسب الإنتروبيا الانتقالية هو الذي يدفع التفاعل.
\end{solution}

\begin{solution}{exo:b3:supramolecular:7}
$\Delta G^\circ = -RT\ln(\num{2.0e5}) = \qty{-30.2}{kJ/mol}$؛ $T\Delta S^\circ = \Delta H^\circ - \Delta G^\circ = -40 + 30.2 = \qty{-9.8}{kJ/mol}$؛
$\Delta S^\circ = \qty{-33}{J.K^{-1}.mol^{-1}}$: ارتباط تدفعه الأنتالبي، يُدفع جزء من ثمنه بالإنتروبيا.
\end{solution}

\begin{solution}{exo:b3:supramolecular:8}
يربط النحاس(I) وحدتي فينانترولين ثنائيتي المخلب، ولأنه رباعي الأوجه، يمسكهما
بزاوية قائمة، إحداهما منظومة عبر الحيّز الذي ستنغلق فيه حلقة
الأخرى. ثم يشبكهما إغلاق السلسلتين كلتيهما في حلقتين. وفائض كبير من
السيانيد، الذي يرتبط بالنحاس(I) أقوى بكثير، ينزع الفلز ويترك
\omterm{def:b3:supramolecular:interlocked}{الكاتينان} الحر.
\end{solution}

\begin{solution}{exo:b3:supramolecular:9}
$k_c = \pi\Delta\nu/\sqrt2 = \pi \times 120/1.414 = \qty{267}{s^{-1}}$. أيرنغ:
$\Delta G^\ddagger = RT\ln(k_BT/hk_c) = 8.314 \times 300 \times \ln(\num{6.25e12}/267) = \qty{59.6}{kJ/mol}$.
\end{solution}

\begin{solution}{exo:b3:supramolecular:10}
مع فائض من الدواء الصلب، يثبت تركيز الدواء الحر عند $S_0$. ثم $[\mathrm{DC}] = KS_0[\mathrm C]$ 
و$C_t = [\mathrm C] + [\mathrm{DC}] = [\mathrm C](1 + KS_0)$، ومنه $[\mathrm{DC}] = KS_0C_t/(1 + KS_0)$ والذوبانية الكلية $S_0 + [\mathrm{DC}]$. وهنا
$KS_0 = 0.050$: $S = 0.10 + 0.050 \times 10/1.050 = \qty{0.58}{mmol/L}$، أي نحو ستة أضعاف الذوبانية الذاتية.
\end{solution}

\begin{solution}{exo:b3:supramolecular:11}
حين $1/K \ll H_0, G_0$، $[\mathrm{HG}] \approx \bigl(H_0 + G_0 - |H_0 - G_0|\bigr)/2 = \min(H_0, G_0)$: فالنسبة المرتبطة هي $G_0/H_0$ حتى
مكافئ واحد، ثم 1. ولا يعود المنحنى يتوقف على $K$: فكل قيمة كبيرة بما يكفي
تعطي الانكسار الحاد نفسه في حدود الضجيج، فلا تبيّن المعطيات إلا أن $K$ يتجاوز
حدًا أدنى ما.
\end{solution}

\begin{solution}{exo:b3:supramolecular:12}
$k_{\mathrm{off}} = k_{\mathrm{on}}/K = \qty{1e9}{L.mol^{-1}.s^{-1}}/\qty{1500}{L.mol^{-1}} = \qty{6.7e5}{s^{-1}}$. وفرق الانزياح
$\qty{0.160}{ppm}$ يساوي \qty{64}{Hz} عند \qty{400}{MHz}؛ ويحتاج الاندماج إلى $k \approx \pi \times 64/\sqrt2 = \qty{140}{s^{-1}}$. فالتبادل
أسرع بآلاف المرات: إشارة واحدة متوسطة.
\end{solution}

\begin{solution}{pb:b3:supramolecular:1}
\textbf{1.} حلقة من ست وحدات $\ce{-CH2CH2O-}$؛ وفي المعقد يجلس $\ce{K+}$ في المركز، مرتبطًا بذرات
الأكسجين الست المتجهة إلى الداخل، ومجموعات $\ce{CH2}$ في الخارج.

\textbf{2.} $\ce{K+}$: $-RT\ln 10 \times 6.1 = \qty{-34.8}{kJ/mol}$؛ $\ce{Na+}$: \qty{-25.1}{kJ/mol}.

\textbf{3.} $10^{6.1 - 4.4} = 50$.

\textbf{4.} الأيون $\ce{K+}$ الأكبر (\qty{138}{pm}) يبلغ ذرات أكسجين الحلقة الست في آن واحد؛ أما
$\ce{Na+}$ الأصغر (\qty{102}{pm}) فلا يستطيع، ما لم تنطوِ الحلقة، وهذا يكلّف إجهادًا.

\textbf{5.} سطحه الخارجي هيدروكربوني وشحنته مدفونة؛ أما $\ce{KCl}$
فيحتاج إلى إذابة أيوناته، وهذا ما لا يقدر عليه التولوين.

\textbf{6.} يذيب الماء $\ce{K+}$ أفضل بكثير من الميثانول، ويجب أن يفقد الأيون ذلك
الماء ليدخل الحلقة.

\textbf{7.} $\ce{K+}(\text{aq}) + \ce{MnO4-}(\text{aq}) + \mathrm L(\text{عضوي}) \rightleftharpoons \ce{[KL]+MnO4-}(\text{عضوي})$.

\textbf{8.} يبقى $[\ce{K+}]_{\mathrm{aq}}$ مساويًا \qty{0.10}{mol/L}: $y = \num{1.0e5} \times 0.10 \times (\num{1.0e-3} - y)^2 = \num{1.0e4}(\num{1.0e-3} - y)^2$.

\textbf{9.} مع $u = \num{1.0e-3} - y$: $\num{1.0e4}u^2 + u - \num{1.0e-3} = 0$، $u = \qty{2.70e-4}{mol/L}$، $y = \qty{7.30e-4}{mol/L}$: أي 73\,\%
مستخلصة.

\textbf{10.} تعطي $y = \num{1.0e4}(\num{1.0e-3} - y)(\num{1.0e-2} - y)$ القيمة $y = \qty{9.89e-4}{mol/L}$: أي 99\,\%.

\textbf{11.} تعطي $y = \num{1.0e4}(\num{1.0e-3} - y)(\num{1.0e-4} - y)$ القيمة $y = \qty{9.0e-5}{mol/L}$: أي 9\,\%، يحدّها \omterm{def:b3:supramolecular:hosts}{الإيثر التاجي}،
المحمَّل بنسبة 90\,\%.

\textbf{12.} يحتفظ أيون البرمنغنات بلونه البنفسجي في التولوين. وهو هناك
غير مذاب، أنيون عارٍ، فيتفاعل أسرع بكثير منه في الماء، حيث
يحجبه غلاف إماهة.

\textbf{13.} يأتي أيون الصوديوم ويذهب أسرع بكثير من الفرق بين
ترددي الرنين: تبادل سريع.

\textbf{14.} $\delta = \delta_{\mathrm{free}} + (\delta_{\mathrm{bound}} - \delta_{\mathrm{free}})[\mathrm{HG}]/H_0$، مع $[\mathrm{HG}]$ من متساوي الحرارة الدقيق، و$H_0 = \qty{2.0e-3}{mol/L}$،
و$G_0 = (\text{المكافئات})\,H_0$، و$\delta_{\mathrm{free}} = \qty{3.560}{ppm}$.

\textbf{15.} $K = (1.55 \pm 0.08) \times 10^3\ \unit{L.mol^{-1}}$, $\delta_{\mathrm{bound}} = \qty{3.719 \pm 0.001}{ppm}$.

\textbf{16.} الجذر التربيعي لمتوسط مربعات البواقي \qty{0.0014}{ppm}، أي بحجم دقة
القراءة، دون أي اتجاه: فالنموذج 1:1 مطابق.

\textbf{17.} النسبة المرتبطة $f$ عند $G_0 = fH_0 + f/(K(1 - f))$: 0.28 مكافئ من أجل $f = 0.2$ و2.1 من أجل $f = 0.8$.

\textbf{18.} $KH_0 = 10^{6.1} \times \num{2.0e-3} \approx 2500$: فالمنحنى سينكسر بحدة عند مكافئ واحد
ولن يعطي إلا حدًا أدنى؛ ويلزم إجراء تجربة تنافس أو قياس حراري عند
تركيز أدنى.

\textbf{19.} في التولوين، في محلول متجانس من الألكين والبرمنغنات.

\textbf{20.} يغادر المنغنيز المختزل الطور العضوي (في صورة ثنائي أكسيد المنغنيز)،
ويلتقط \omterm{def:b3:supramolecular:hosts}{الإيثر التاجي}، المتحرر عند السطح الفاصل، أيونًا آخر من $\ce{K+}$ و$\ce{MnO4-}$. فكل جزيء
من \omterm{def:b3:supramolecular:hosts}{الإيثر التاجي} ينقل البرمنغنات مرات كثيرة: فتكفي كمية حفزية.

\textbf{21.} يضرب $K_{\mathrm{ex}}$ في $[\ce{K+}]$: فتركيز البوتاسيوم المرتفع يدفع الاستخلاص.

\textbf{22.} أملاح الأمونيوم الرباعية، مثل كلوريد رباعي بوتيل الأمونيوم أو
كلوريد بنزيل ثلاثي إيثيل الأمونيوم، التي يكون كاتيونها محبًا للدهون بذاته.

\textbf{23.} يسبب البنزين ابيضاض الدم؛ ويؤدي التولوين العمل نفسه بخطر أصغر
بكثير.

\textbf{24.} مع \qty{1.0}{mmol/L} من \omterm{def:b3:supramolecular:hosts}{الإيثر التاجي}، يُستخلص 73\,\% من البرمنغنات إلى
التولوين.
\end{solution}
